\section*{Thesen}

\textbf{Projekt-Thesis}

\textbf{Entwicklung einer mobilen Begleit-Applikation zur datenbankgestützten Verwaltung von Hackintosh-Systemen mit Flutter und SQLite}

\vspace{0.5cm}

\textbf{Eingereicht am:} 11.04.2026

\textbf{von:} Kilian Ketelhohn (520858)

\textbf{geboren am:} 27.11.2001 in Schwerin

\textbf{Betreuer:} Prof.\ Dr.\ Matthias Kreuseler

\vspace{1cm}

\begin{itemize}

\item Die erfolgreiche Konfiguration von Hackintosh-Systemen erfordert eine strukturierte Erfassung technischer Parameter, die insbesondere in Phasen ohne Netzanbindung (Installation) mobil verfügbar sein müssen.

\item Eine relationale Datenbank wie SQLite bietet durch die Unterstützung von Fremdschlüssel-Beziehungen und ACID-Garantien die notwendige Integrität, um komplexe Abhängigkeiten zwischen Hardware und Treibern (Kexts) fehlerfrei abzubilden.

\item Das Cross-Plattform-Framework Flutter ermöglicht durch seine eigene Rendering-Engine eine performante Implementierung von technisch anspruchsvollen Features wie 3D-Hardware-Visualisierungen und Echtzeit-Boot-Animationen auf mobilen Endgeräten.

\item Ein Offline-First-Ansatz in Verbindung mit einer REST-API-Synchronisation stellt sicher, dass Systemdaten lokal redundant vorliegen und dennoch mit einer zentralen Community-Datenbank abgeglichen werden können.

\item Die Nutzung von gerätespezifischer Sensorik, insbesondere das Scannen von QR-Codes, minimiert die Fehlerquote bei der manuellen Übertragung komplexer Konfigurations-Strings zwischen verschiedenen Endgeräten.

\item Die gezielte Optimierung der Grafik-Pipeline (Wechsel auf den Skia-Renderer) ermöglicht eine flüssige Ausführung moderner App-Architekturen auch auf ressourcenlimitierter Legacy-Hardware wie dem Samsung Galaxy A3 (2016).

\item Die Kombination aus interaktiver 3D-Grafik und strukturierter Datenbank-Abfrage erhöht die Effizienz bei der Fehlerdiagnose und Hardware-Identifikation im Hackintosh-Bereich signifikant.

\end{itemize}